%% =====================================================================
%%  B1-T-13-02 -- what B1 contributes to "Weakness as a Mask"
%%  -------------------------------------------------------------------
%%  LAYER A: APPEND-ONLY. Written once at the entry's write, then left
%%  alone. If the source has more to say later, add a bullet -- do not
%%  rewrite. Material found ONLY here goes in the `unique` slot with
%%  @unique set to yes, which makes it undeletable (rule R10).
%% =====================================================================
%% @kind: source
%% @id: B1-T-13-02
%% @book: B1
%% @topic: T-13-02
%% @locator: Tactic 5 (ch. VI), pp. 59--70
%% @status: distilled
%% @unique: yes
%% @integrated: yes
%% @updated: 2026-09-19
% @allow-rewrite: skeleton filled at the write of T-13-02 (planned -> stable lifecycle)

\dossier{B1}{T-13-02}{\bookshort{B1}}{Tactic 5 (ch. VI), pp. 59--70}

\begin{distilled}
  \item The meek-house analogy: the meek inherit sales because nobody bothers to defend against them.
  \item Never know much and you get further: the salesman who out-earned mocking tradesmen underplayed knowledge and sharpness.
  \item A man who believes he is the sharp one stops checking and can be taken for more money.
  \item Displayed ignorance is the grant that lets the counterpart teach, disclose, and overextend while feeling superior.
\end{distilled}
\begin{unique}
  \item The specific instruction to underplay knowledge in house-to-house selling, priced by the author as worth more than the tradesmen's jeers he absorbed.
\end{unique}
\begin{terms}
  \item \emph{the meek cover} --- Sparkman's term for the deliberate under-display of knowledge and status that lowers the counterpart's guard.
\end{terms}
